\addcontentsline{toc}{subsubsection}{Solution 30. The Quadratic--Difference Squeeze}

\begin{dlsummary}[title={\bfseries Solution 30. The Quadratic--Difference Squeeze},breakable]

\begin{proof}
Taking $y=x$ in the first inequality gives $f(x)\geq x^2$ for every
real $x$. Define
\[
g(x)=f(x)-x^2.
\]
Then $g(x)\geq0$. Substituting $f(t)=t^2+g(t)$ into the second
inequality and cancelling the quadratic terms yields
\[
g(x+y)+g(x-y)\leq2g(x)\qquad(x,y\in\mathbb R). \tag{1}
\]

\medskip
\noindent
\textbf{Step 1. Show that $g$ is constant.}

Fix arbitrary $a,b\in\mathbb R$, and set
\[
x_k=a+k(b-a)\qquad(k\in\mathbb Z),
\qquad
\Delta_k=g(x_{k+1})-g(x_k).
\]
Applying (1) with $x=x_k$ and $y=b-a$ gives
\[
g(x_{k+1})+g(x_{k-1})\leq2g(x_k),
\]
so $\Delta_k\leq\Delta_{k-1}$. In particular,
$\Delta_j\leq\Delta_0=g(b)-g(a)$ for every integer $j\geq0$.
For any positive integer $n$, summing these bounds gives
\[
-g(a)\leq g(x_n)-g(a)
=\sum_{j=0}^{n-1}\Delta_j
\leq n\bigl(g(b)-g(a)\bigr),
\]
where the first inequality uses $g(x_n)\geq0$.
If $g(b)<g(a)$, choosing $n$ sufficiently large makes the rightmost
quantity less than $-g(a)$, a contradiction. Thus $g(b)\geq g(a)$.
Since $a,b$ were arbitrary, interchanging them shows $g(a)=g(b)$;
therefore $g$ is constant. Write $g(x)=d$. Its nonnegativity gives
$d\geq0$, and hence
\[
f(x)=x^2+d\qquad(x\in\mathbb R).
\]

\medskip
\noindent
\textbf{Verification.}

For $d\geq0$,
\[
f(x)+f(y)-2xy=(x-y)^2+2d\geq0,
\]
and
\[
f(x+y)+f(x-y)
=(x+y)^2+(x-y)^2+2d
=2f(x)+2y^2.
\]
Thus every function in the stated family satisfies both inequalities.
\end{proof}

\end{dlsummary}